% Generated shared reviewed material.
\section{Requested-error theta samples for three fixed conductors}\label{at:sampling}
For the primitive characters $\chi_3=(0,1,-1)$, $\chi_4=(0,1,0,-1)$,
$\chi_5=(0,1,-1,-1,1)$ with parities $a=1,1,0$, retain the classical
primitive theta/Mellin identity and root number one already stated above.
Put $D=[-1,2]+i[-4/3,4/3]$ and $\delta=3/q$.
The following adaptive extension is implemented with independent interval
and sample checks; it is an informal analytic theorem, not a Lean result.
\begin{theorem}[Finite adaptive sample specification]\label{at:sample}
For every supported conductor, exact rational $s\in D$, and rational
$\eta>0$, the choices below terminate and give a rational complex midpoint
$y$ and exact rational radius $R\le\eta$ with
$|y-\Lambda_q(s)|\le R$.
The same conclusion holds for a specified true circle node when its
root and node enclosures are refined on every precision retry.
\end{theorem}
\begin{proof}
For $x\ge0$ define the exact rational bound $E(x)=2^{-\lfloor x\rfloor}$.
Since $e>2$, $e^{-x}\le E(x)$; this bound is nonincreasing and tends to zero.
Let $Q(x)=x^{-1}+2x^{-2}+2x^{-3}$ and
$C_0=1$, $C_1=(1+\delta^{-1})/2$.
For $N\ge1$, $n_0=N+1$, $J\ge1$, $T=2^J$, use
\begin{align*}
 B_n&=2C_aQ(\delta n_0^2)
       \frac{E(\delta n_0^2/2)}{1-E(\delta n_0)},\\
 B_t&=\sum_{n=1}^N2n^aT^3E(\delta n^2T)Q(\delta n^2T).
\end{align*}
Here $\delta n_0\ge6/5$ ensures a positive denominator; $E(x)$ equals
one for $0\le x<1$, so this safeguard is essential.
On the four panels of each $[2^j,2^{j+1}]$, write midpoint $m$,
halfwidth $h$, $\rho=m/2$, $v=h/\rho\le2/9$,
and $B(m)=\max\{(m+\rho)^2,(m-\rho)^{-2},1\}$.
Use the Cauchy remainder
\[
 B_Q=\sum_{\mathrm{panels}}\sum_{n=1}^N
 \frac{2n^a\,2h\,3B(m)E(\delta n^2m/2)v^{2K+2}}
 {(2K+3)(1-v^2)}.
\]
Choose in turn the first $N\ge1$, first $J\ge1$, and first $K\ge0$
for which their respective budgets are at most $\eta/8$.
Gaussian decay, then exponential time decay, then $v\le2/9$ prove
termination. These are the accepted tail and Cauchy arguments with each
true exponential separately replaced by its upper bound $E$.

There are $P=4JN$ potential pairs. A nonzero-character pair may be
omitted only if its entire two-power contribution has bound
\[
 b=4hn^a(m+h)^2E(\delta n^2(m-h))\le\eta/(8P).
\]
Add every omitted bound to $B_{\rm discard}$, giving
$B_{\rm discard}\le\eta/8$.
Evaluate all other integrated Taylor polynomials by the stated exact
coefficient recurrence using outward rectangles. Refine both interval
precision and all primitive series remainders until the final rectangle's
sum of coordinate halfwidths satisfies $R_{\rm arith}\le\eta/4$.
For this fixed finite expression all input and primitive intervals shrink;
continuity and outward arithmetic imply $R_{\rm arith}\to0$.
Precision doubling therefore terminates, and the sum of the five raw
components is at most $3\eta/4$.
Choose an integer output scale $B$ such that $5\,2^{-B}\le\eta/8$ and
round each component upward by
$\operatorname{up}_B(x)=2^{-B}\lceil2^Bx\rceil$.
The reported radius is their literal exact sum. With
$\mathbf b=(R_{\rm arith},B_n,B_t,B_Q,B_{\rm discard})$, it obeys
\[
 |y-\Lambda_q(s)|\le\sum_{i=1}^5\operatorname{up}_B(b_i)
 \le\frac{3\eta}{4}+5\,2^{-B}\le\frac{7\eta}{8}<\eta.
\]
The output retains exact rational midpoints, rounded components, $B$,
raw interval radius, parameters, input metadata and source hashes.
The raw analytic bounds are exactly reconstructible; they need not be
serialized as potentially enormous fractions. The output precision varies
with the request and introduces no fixed reporting floor.
\end{proof}

This fixed-conductor sample result uses no hypothesis on zeros.
It does not remove a downstream fixed-node Fourier alias floor, supply
an all-conductor theorem or a positive arithmetic determinant, or prove GRH.
